% !TeX root = ../../../金融经济学_王江_讲义及习题.tex
% 取 A=1.35、B=0.88，曲线和焦点使用同一组参数。
\begin{tikzpicture}[x=.84cm,y=.84cm,font=\small,>=Stealth]
 \def\wjHyperbolaA{1.35}
 \def\wjHyperbolaB{.88}
 \pgfmathsetmacro{\wjHyperbolaC}{sqrt(\wjHyperbolaA^2+\wjHyperbolaB^2)}
 \pgfmathsetmacro{\wjHyperbolaM}{\wjHyperbolaB/\wjHyperbolaA}

 \coordinate (F1) at (-\wjHyperbolaC,0);
 \coordinate (F2) at (\wjHyperbolaC,0);
 \coordinate (V1) at (-\wjHyperbolaA,0);
 \coordinate (V2) at (\wjHyperbolaA,0);
 \coordinate (P) at ({\wjHyperbolaA*sqrt(1+(1.34/\wjHyperbolaB)^2)},1.34);

 \draw[gray,densely dashed] (-3.23,{-3.23*\wjHyperbolaM}) --
   (3.23,{3.23*\wjHyperbolaM});
 \draw[gray,densely dashed] (-3.23,{3.23*\wjHyperbolaM}) --
   (3.23,{-3.23*\wjHyperbolaM});
 \draw[->] (-3.38,0) -- (3.47,0) node[right] {$x$};
 \draw[->] (0,-2.13) -- (0,2.23) node[above] {$y$};
 \node[below left,inner sep=2pt] at (0,0) {$O$};

 \draw[teal!75!black,thin] (F1) -- (P) -- (F2);
 \foreach \sgn in {-1,1} {
   \draw[blue!70!black,thick,domain=-1.93:1.93,samples=81,variable=\t]
     plot ({\sgn*\wjHyperbolaA*sqrt(1+(\t/\wjHyperbolaB)^2)},\t);
 }

 \fill[teal!75!black] (F1) circle (1.65pt) node[below left,inner sep=2pt] {$F_1$};
 \fill[teal!75!black] (F2) circle (1.65pt) node[below right,inner sep=2pt] {$F_2$};
 \fill[blue!70!black] (V1) circle (1.65pt) node[above right,inner sep=2pt] {$V_1$};
 \fill[blue!70!black] (V2) circle (1.65pt) node[above left,inner sep=2pt] {$V_2$};
 \fill[teal!75!black] (P) circle (1.65pt) node[right,inner sep=3pt] {$P$};
\end{tikzpicture}
